\section{Code Structure and Solution Layout}
\label{sec:impl-structure}
% Authoring brief for this section lives in section.md (§6.1).

\Cref{sec:design-style,sec:design-modules} described the architecture as concerns, ports, and a rule that dependencies point inward.
This chapter goes inside that description to show how it is realised in code, beginning here with where the code lives before \cref{sec:impl-session} follows a session through it.
The system is a monorepo with two independent roots, a .NET backend solution under \texttt{Backend/} and a Next.js front-end under \texttt{Frontend/}, with no shared workspace manifest binding them, so that each is built and versioned on its own.
Their folder layout is shown in \cref{fig:impl-tree} and their dependency structure in \cref{fig:impl-structure}.

\begin{figure}[htbp]
\footnotesize
\dirtree{%
.1 \texttt{reading-the-reader-monorepo/}.
.2 \texttt{Backend/src/} \DTcomment{.NET solution}.
.3 \texttt{core/}.
.4 \texttt{ReadingTheReader.core.Domain/} \DTcomment{domain entities}.
.4 \texttt{ReadingTheReader.core.Application/} \DTcomment{ports, services, snapshots}.
.5 \texttt{ApplicationContracts/Realtime/} \DTcomment{the four concern modules}.
.3 \texttt{infrastructure/}.
.4 \texttt{ReadingTheReader.TobiiEyetracker/} \DTcomment{Tobii adapter}.
.4 \texttt{ReadingTheReader.RealtimeMessenger/} \DTcomment{WebSocket broadcaster}.
.4 \texttt{ReadingTheReader.Realtime.Persistence/} \DTcomment{file store + checkpoints}.
.3 \texttt{ReadingTheReader.WebApi/} \DTcomment{transport, composition root}.
.2 \texttt{Frontend/src/} \DTcomment{Next.js application}.
.3 \texttt{app/} \DTcomment{routes only (sidebar groups)}.
.3 \texttt{modules/pages/} \DTcomment{feature UI}.
.3 \texttt{redux/}, \texttt{lib/} \DTcomment{state, API, gaze socket}.
.3 \texttt{components/}, \texttt{hooks/} \DTcomment{shared UI}.
}
\caption{Curated layout of the monorepo (selected folders only). The backend groups projects into \texttt{core} (the \texttt{Domain} entities and the \texttt{Application} with its ports), \texttt{infrastructure} (the adapters), and the \texttt{WebApi} composition root; the four concern modules of \cref{sec:design-modules} appear as folders beneath the application contracts. The front-end separates routing (\texttt{app/}) from feature UI, shared state and transport, and shared components.}
\label{fig:impl-tree}
\end{figure}

\begin{figure}[htbp]
\centering
\resizebox{\linewidth}{!}{%
\begin{tikzpicture}[
  web/.style={rectangle, rounded corners=3pt, draw=dtured, very thick, fill=dtured!10, align=center, text width=2.7cm, minimum height=0.95cm, inner sep=3pt, font=\scriptsize},
  infra/.style={rectangle, rounded corners=3pt, draw=black!55, thick, fill=grey!25, align=center, text width=2.2cm, minimum height=0.95cm, inner sep=3pt, font=\scriptsize},
  app/.style={rectangle, rounded corners=3pt, draw=navyblue, very thick, fill=blue!10, align=center, text width=2.6cm, minimum height=0.95cm, inner sep=3pt, font=\scriptsize},
  dom/.style={rectangle, rounded corners=3pt, draw=navyblue, very thick, fill=blue!18, align=center, text width=2.0cm, minimum height=0.8cm, inner sep=3pt, font=\scriptsize},
  fe/.style={rectangle, rounded corners=3pt, draw=navyblue, thick, fill=navyblue!8, align=center, text width=3.4cm, minimum height=1.0cm, inner sep=3pt, font=\scriptsize},
  dep/.style={->, thick, black!60},
  fedep/.style={->, semithick, black!45},
  ttl/.style={font=\small\bfseries, text=black!75},
  note/.style={font=\scriptsize\itshape, text=black!70, align=center},
]
% ---- panel titles ----
\node[ttl] at (-2.3, 6.3) {Backend solution (.NET)};
\node[ttl] at ( 5.0, 6.3) {Frontend (Next.js)};
% ---- divider ----
\draw[dashed, black!30] (2.9,-1.0) -- (2.9,6.0);
% ---- backend projects ----
\node[web]   (web) at (-4.7, 5.3) {{\bfseries WebApi}\\{\scriptsize transport + composition root}};
\node[infra] (per) at (-4.7, 3.5) {{\bfseries Realtime}\\{\scriptsize Persistence}};
\node[infra] (msg) at (-1.7, 3.5) {{\bfseries Realtime}\\{\scriptsize Messenger}};
\node[infra] (tob) at ( 1.4, 3.5) {{\bfseries Tobii}\\{\scriptsize Eyetracker}};
\node[app]   (app) at (-1.7, 1.7) {{\bfseries Application}\\{\scriptsize contracts, orchestration}};
\node[dom]   (dom) at (-1.7, 0.0) {{\bfseries Domain}};
% ---- backend references (A -> B: A depends on B) ----
\draw[dep] (app) -- (dom);
\draw[dep] (per) -- (app);
\draw[dep] (msg) -- (app);
\draw[dep] (tob) -- (app);
\draw[dep] (tob) -- (dom);
\draw[dep] (web) -- (per);
\draw[dep] (web) -- (msg);
\draw[dep] (web) -- (tob);
\draw[dep] (web) -- (app);
% ---- frontend layer stack ----
\node[fe] (feapp)   at (5.0, 5.1) {{\bfseries app/}\\{\scriptsize routes \& layouts (with-/without-sidebar)}};
\node[fe] (femod)   at (5.0, 3.6) {{\bfseries modules/pages}\\{\scriptsize feature UI}};
\node[fe] (festate) at (5.0, 2.1) {{\bfseries redux} \& {\bfseries lib}\\{\scriptsize state, RTK Query API, gaze socket}};
\node[fe] (feshare) at (5.0, 0.6) {{\bfseries components} \& {\bfseries hooks}\\{\scriptsize shared UI}};
\draw[fedep] (feapp)   -- (femod);
\draw[fedep] (femod)   -- (festate);
\draw[fedep] (festate) -- (feshare);
% ---- legend ----
\node[note, anchor=north west] at (-6.1,-1.2) {Arrow $A \rightarrow B$: project (or layer) $A$ references / depends on $B$. In the backend every reference points inward, towards the \textbf{Domain} core.};
\end{tikzpicture}%
}
\caption{Code structure of the two monorepo roots. In the backend solution every project reference points inward (\texttt{WebApi}, the composition root, depends on \texttt{Application} and all infrastructure; each infrastructure adapter and \texttt{Application} depend inward towards \texttt{Domain}), which is the dependency rule of \cref{sec:design-style} enforced by the build rather than merely intended. The front-end is a single Next.js application whose source separates thin routing (\texttt{app/}) from feature UI (\texttt{modules/pages}), shared state and transport (\texttt{redux}, \texttt{lib}), and shared components.}
\label{fig:impl-structure}
\end{figure}

The backend solution realises the concentric structure of \cref{sec:design-style} directly as projects.
A \texttt{Domain} project at the centre holds the basic entities, an \texttt{Application} project holds the use-case services, the snapshots, and the port interfaces, three infrastructure projects supply the adapters for the Tobii device, realtime messaging, and persistence, and a \texttt{WebApi} project provides the HTTP and WebSocket transport and composes the whole at startup.
Each conceptual layer of the architecture is thus a separate compilation unit, which is what allows the dependency rule to be enforced rather than merely intended.

That rule is visible in the project references, and because a project cannot compile against a reference it does not declare, the references are evidence rather than documentation.
The \texttt{Domain} project references nothing; the \texttt{Application} project references only \texttt{Domain}; and each infrastructure adapter references inward to \texttt{Application}, with the Tobii project additionally referencing \texttt{Domain}, so that every reference points towards the centre, as \cref{fig:impl-structure} shows.
The one project that references outward is \texttt{WebApi}, which depends on \texttt{Application} and on all three infrastructure projects at once.
This is deliberate, and it is the single exception the style permits: as the composition root, \texttt{WebApi} alone must know the concrete adapters in order to register them against the core's ports at startup, while the core itself never names them.

The front-end is a single Next.js application whose source mirrors the same discipline of separating routing from logic.
The \texttt{app/} directory contains only routes and layouts, organised into two route groups, a researcher-facing group with the operator sidebar and a participant-facing group without it, which are the two surfaces of the container view of \cref{sec:design-style}.
Feature UI lives in \texttt{modules/pages}, cross-feature state and the typed REST clients in \texttt{redux}, and shared helpers including the realtime gaze socket in \texttt{lib}, with common components and hooks factored into \texttt{components} and \texttt{hooks}.
Routing therefore stays thin and hands off to feature logic that is kept separately, in the same spirit as the backend's separation of transport from core.

Because there is no generated package shared between the two roots, the TypeScript types the front-end uses are hand-written mirrors of the backend records, and keeping the two aligned is a manual discipline rather than a compiler-checked one.
With the territory mapped, the next section puts it in motion: \cref{sec:impl-session} follows a single experiment session through these projects, from the researcher's first setup action to the sealed record left behind.
